\section{Restriction and Matched Multiresolution Geometry}

\subsection{Point-current restriction}

Point interpolation and raster analysis involve different restriction
functionals. Let \(d_h\) be the forward difference on an interpolated line,
and let \(B_s\) sum each block of \(s\) consecutive fine differences. Using
cardinality gives
\begin{align}
 (B_s d_h\Tcal_s y)_i
 &=\sum_{r=0}^{s-1}
 \bigl[(\Tcal_s y)_{is+r+1}-(\Tcal_s y)_{is+r}\bigr]\nonumber\\
 &=y_{i+1}-y_i.
 \label{eq:telescoping}
\end{align}
Consequently,
\begin{equation}
 D_s\Tcal_s=I,\qquad B_s d_h\Tcal_s=d_c,
 \label{eq:restriction}
\end{equation}
where \(D_s\) selects source sites and \(d_c\) is the source forward
difference. Thus \eqref{eq:restriction} identifies the conserved mass of the
interpolating current; site selection itself is not used as an antialiased
raster downscale.

\subsection{Conservative scale coordinates}

The centered-moment construction below uses the fixed three-point
Gauss--Legendre nodes and weights
\begin{equation}
 \bm\xi=(-\sqrt{3/5},0,\sqrt{3/5}),\qquad
 \bm w=(5/9,8/9,5/9).
 \label{eq:gauss3}
\end{equation}

At a multiresolution interface, raster entries are interpreted as cell
averages. Accordingly, for an even fine line
\(x=(x_0,\ldots,x_{2N-1})\), define the parent average and centered child
moment by
\begin{align}
 z_i&=\frac{x_{2i}+x_{2i+1}}{2},
 \label{eq:scaleaverage}\\
 \mu_i^\star&=\frac{x_{2i+1}-x_{2i}}{2}.
 \label{eq:actualmoment}
\end{align}
The resulting analysis low-pass is positive and normalized, with transfer
function
\begin{equation}
 H(\omega)=\frac{1+e^{-\mathrm i\omega}}{2}
 =e^{-\mathrm i\omega/2}\cos(\omega/2),
 \label{eq:nyquistresponse}
\end{equation}
which yields \(H(\pi)=0\) exactly.

The centered-moment proposal is obtained from the already defined CONV
profile, rather than from a second interpolator. Using the Gauss--Legendre
nodes and weights in \eqref{eq:gauss3}, for \(1\leq i<N-1\), put
\begin{align}
 \bar z_i^-&=\frac12\sum_{k=0}^{2}w_k
 (\Tcal z)(i-\tfrac14+\tfrac14\xi_k),\nonumber\\
 \bar z_i^+&=\frac12\sum_{k=0}^{2}w_k
 (\Tcal z)(i+\tfrac14+\tfrac14\xi_k),
 \label{eq:halfcells}\\
 p_i&=\frac{\bar z_i^+-\bar z_i^-}{2}.
 \label{eq:momentproposal}
\end{align}
Since \(\Tcal z\) is quintic on each half-cell, three-point Gauss--Legendre
quadrature is exact. At the two outer cells, the fixed closure is
\begin{equation}
 p_0=m_0/4,\qquad p_{N-1}=m_{N-1}/4,
 \label{eq:momentboundary}
\end{equation}
using the one-sided jet in \eqref{eq:mbnd}; this proposal is exact for affine
cell averages, for which \(p_i=(z_{i+1}-z_i)/4\).

\subsection{Admissible moment polytope}

Put \(\Delta_i=z_{i+1}-z_i\), fill zero signs by the ledger rule of
Section~V, and denote the resulting interval signs by \(s_i\). When every
\(\Delta_i\) vanishes, set \(\Pcal(z,t)=\{0\}\) and
\(\widehat\mu=0\). Otherwise the filled signs are nonzero, and each cell is
assigned a sign \(t_i\), with \(t_0=s_0\),
\(t_{N-1}=s_{N-2}\), and, in the interior,
\begin{equation}
 t_i=\begin{cases}
 s_{i-1},&s_{i-1}=s_i,\\
 \arg\max_{s\in\{s_{i-1},s_i\}}sp_i,&s_{i-1}\ne s_i,
 \end{cases}
 \label{eq:momentsign}
\end{equation}
with the left sign selected on a tie. The second line therefore locates an
existing turning point in one of its two adjacent child cells.

When \(\mu\) is synthesized with zero detail, the fine currents are
\begin{equation}
 g_{2i}=2\mu_i,\qquad
 g_{2i+1}=\Delta_i-\mu_i-\mu_{i+1}.
 \label{eq:finecurrentmoment}
\end{equation}
Consequently, for fixed \(t\), the complete admissible set is
\begin{equation}
 \Pcal(z,t)=\left\{\mu:
 t_i\mu_i\geq0,\quad
 s_i(\Delta_i-\mu_i-\mu_{i+1})\geq0
 \right\}.
 \label{eq:momentpolytope}
\end{equation}
\begin{lemma}[Compact moment fibre]
For every nonconstant \(z\) and every assignment \(t\) produced by
\eqref{eq:momentsign}, \(\Pcal(z,t)\) is a nonempty compact convex polytope.
\end{lemma}

\begin{proof}
The set is a finite intersection of closed half-spaces and contains zero, so
only an unbounded ray remains to be excluded. Let \(h\) belong to its
recession cone and put \(v_i=t_ih_i\geq0\). The homogeneous face inequalities
are
\begin{equation}
 (s_it_i)v_i+(s_it_{i+1})v_{i+1}\leq0.
 \label{eq:momentrecession}
\end{equation}
Suppose \(v_k>0\) for the least such \(k\). If \(k=0\), then \(t_k=s_k\).
If \(k=N-1\), the preceding inequality contradicts \(v_k>0\) because
\(t_{N-1}=s_{N-2}\). If \(0<k<N-1\), the preceding inequality and
\(v_{k-1}=0\) either contradict \(v_k>0\), or force
\(t_k=-s_{k-1}=s_k\). Thus every surviving case has \(k<N-1\) and
\(t_k=s_k\). Equation~\eqref{eq:momentrecession} on face \(k\) then forces
\(t_{k+1}=-s_k\) and \(v_{k+1}\geq v_k\). By
\eqref{eq:momentsign}, this is possible only when
\(s_{k+1}=-s_k\) and \(t_{k+1}=s_{k+1}\). Repetition propagates a positive
nondecreasing sequence to the right endpoint, where
\(t_{N-1}=s_{N-2}\) contradicts the required opposite sign. Hence the
recession cone is \(\{0\}\), and the closed polyhedron is bounded.
\end{proof}

On an increasing run, the constraints reduce to
\begin{equation}
 \mu_i\geq0,\qquad \mu_i+\mu_{i+1}\leq\Delta_i.
 \label{eq:monotonemoment}
\end{equation}

The finite admission contains no limiter constant. First form
\begin{equation}
p_i^+=t_i\max(t_ip_i,0).
 \label{eq:momentorthant}
\end{equation}
Let \(\rho(i)\) denote the maximal sign lineage assigned to cell \(i\), and
let \(\rho(\ell)\) denote the lineage of face \(\ell\). Its proposed
consumption is
\begin{equation}
 C_\ell=\sum_{\substack{j\in\{\ell,\ell+1\}\\
                        \rho(j)=\rho(\ell)}}|p_j^+|.
 \label{eq:momentconsumption}
\end{equation}
and, for each lineage, retain the scale
\begin{equation}
\alpha_\rho=\min\left(
 1,\min_{\substack{\ell:\rho(\ell)=\rho\\C_\ell>0}}
 \frac{|\Delta_\ell|}{C_\ell}\right).
 \label{eq:momentrayscale}
\end{equation}
where the inner minimum is \(+\infty\) when the lineage has no positive
consumption. Oppositely oriented moments at a transition relax rather than
consume the face mass and therefore do not enter \(C_\ell\). The result is the
largest independent lineage-ray scale satisfying every incident capacity
bound, and it is denoted by
\begin{equation}
 \widehat\mu_i=\alpha_{\rho(i)}p_i^+=(\Acal_zp)_i.
 \label{eq:momentadmission}
\end{equation}

\begin{lemma}[Admitted-moment bound]
The admission \eqref{eq:momentadmission} satisfies
\begin{equation}
 \norm{\widehat\mu(z)}_\infty\leq
 \max_\ell|\Delta_\ell|\leq2\norm z_\infty.
 \label{eq:admittedmomentbound}
\end{equation}
\end{lemma}

\begin{proof}
Every cell \(i\) has an incident face \(\ell\) in its assigned lineage, as
forced at the endpoints and by the definition of \(t_i\) in the interior.
Consequently \(C_\ell\geq|p_i^+|\). If \(C_\ell=0\), then
\(\widehat\mu_i=0\). Otherwise
\eqref{eq:momentrayscale} gives
\[
 |\widehat\mu_i|
 =\alpha_{\rho(i)}|p_i^+|
 \leq|\Delta_\ell|.
\]
Finally, \(|z_{\ell+1}-z_\ell|\leq2\norm z_\infty\), which proves the second
inequality.
\end{proof}

\subsection{Exact conservative lifting pair}

The detail coordinate is obtained as the child moment not represented by the
transported coarse variation,
\begin{equation}
 r_i=\mu_i^\star-\widehat\mu_i(z).
 \label{eq:analysisdetail}
\end{equation}
so analysis is \(\mathcal Lx=(z,r)\), while synthesis is
\begin{align}
 x_{2i}&=z_i-\widehat\mu_i(z)-r_i,
 \label{eq:synthleft}\\
 x_{2i+1}&=z_i+\widehat\mu_i(z)+r_i.
 \label{eq:synthright}
\end{align}

\begin{theorem}[Conservative CONV lifting]
For every fine line of length \(2N\), \(N\geq5\), in exact arithmetic,
\begin{equation}
 \mathcal L^{-1}\mathcal Lx=x,
 \qquad R\mathcal L^{-1}(z,r)=z,
 \label{eq:fullcycle}
\end{equation}
where \(R\) is the positive restriction \eqref{eq:scaleaverage}. If \(r=0\),
the synthesized current has no more nonzero sign changes than \(Dz\).
\end{theorem}

\begin{proof}
Substituting \eqref{eq:analysisdetail} into
\eqref{eq:synthleft}--\eqref{eq:synthright} returns
\(z_i\mp\mu_i^\star=x_{2i},x_{2i+1}\). Their average is \(z_i\), proving
both identities. By \eqref{eq:momentpolytope}, the nonzero fine-current signs
are a subsequence of
\[
 t_0,s_0,t_1,s_1,\ldots,s_{N-2},t_{N-1}.
\]
Within a sign run the entries agree. At a transition, \(t_i\) equals one of
its two adjacent signs; that triple therefore contains one transition, not
two. Deleting zero currents cannot increase the count.
\end{proof}

The transform consequently distinguishes exact recovery from
discarded-detail reconstruction. Retaining \(r\) recovers every fine cell,
whereas setting \(r=0\) retains cell mass and transported variation while
omitting the unpredicted centered moment. For two detail states on the same
coarse raster, one obtains
\begin{equation}
 \norm{\mathcal L^{-1}(z,r)-\mathcal L^{-1}(z,r')}_\infty
 =\norm{r-r'}_\infty.
 \label{eq:detailstability}
\end{equation}

\begin{proposition}[Multilevel max-norm stability]
The positive restriction satisfies \(\norm R_\infty\leq1\). A \(J\)-level
synthesis obeys
\begin{equation}
 \norm x_\infty\leq3^J\norm{z^{(J)}}_\infty
 +\sum_{j=1}^{J}3^{j-1}\norm{r^{(j)}}_\infty.
 \label{eq:multilevelstability}
\end{equation}
The constant depends on the declared depth, not the line length.
\end{proposition}

\begin{proof}
Since the two restriction weights are positive and sum to one,
\(\norm R_\infty\leq1\). At one synthesis level, using
\eqref{eq:synthleft}--\eqref{eq:synthright} together with
\eqref{eq:admittedmomentbound} gives
\[
 \norm{x^{(j-1)}}_\infty
 \leq3\norm{x^{(j)}}_\infty+\norm{r^{(j)}}_\infty.
\]
Iterating this recurrence from \(x^{(J)}=z^{(J)}\) yields
\eqref{eq:multilevelstability}; induction on the one-level theorem likewise
shows that ordered sign variation is nonincreasing at every zero-detail
level.
\end{proof}

\subsection{Two-dimensional moment geometry}

For a fine \(2\times2\) block, introduce one invariant mean and three
centered moments,
\begin{align}
 z&=\tfrac14(x_{00}+x_{01}+x_{10}+x_{11}),\nonumber\\
 \mu_x&=\tfrac14(-x_{00}+x_{01}-x_{10}+x_{11}),\nonumber\\
 \mu_y&=\tfrac14(-x_{00}-x_{01}+x_{10}+x_{11}),\nonumber\\
 \mu_{xy}&=\tfrac14(x_{00}-x_{01}-x_{10}+x_{11}).
 \label{eq:blockmoments}
\end{align}
from which the inverse is obtained as
\begin{align}
 x_{00}&=z-\mu_x-\mu_y+\mu_{xy},&
 x_{01}&=z+\mu_x-\mu_y-\mu_{xy},\nonumber\\
 x_{10}&=z-\mu_x+\mu_y-\mu_{xy},&
 x_{11}&=z+\mu_x+\mu_y+\mu_{xy}.
 \label{eq:blockinverse}
\end{align}
Thus the four child cells are replaced without redundancy. The horizontal
and vertical proposals are obtained by applying \eqref{eq:halfcells} along
their corresponding factors, and applying the same half-cell functional to
the horizontal proposal along the vertical factor yields the mixed proposal.

Inside a block, the horizontal face currents are
\begin{equation}
 2(\mu_x-\mu_{xy}),\qquad2(\mu_x+\mu_{xy}),
 \label{eq:internalhorizontal}
\end{equation}
and the vertical face currents are
\begin{equation}
 2(\mu_y-\mu_{xy}),\qquad2(\mu_y+\mu_{xy}).
 \label{eq:internalvertical}
\end{equation}
For horizontally adjacent blocks \(A,B\), put \(\Delta_x=z_B-z_A\). Their
upper and lower cross-face currents are
\begin{align}
 h^-_{AB}&=\Delta_x-\mu_{x,A}-\mu_{x,B}
 +\mu_{y,A}-\mu_{y,B}+\mu_{xy,A}+\mu_{xy,B},\nonumber\\
 h^+_{AB}&=\Delta_x-\mu_{x,A}-\mu_{x,B}
 -\mu_{y,A}+\mu_{y,B}-\mu_{xy,A}-\mu_{xy,B}.
 \label{eq:horizontalfaces}
\end{align}
For vertically adjacent blocks \(A,C\), put \(\Delta_y=z_C-z_A\). Their
left and right currents are
\begin{align}
 v^-_{AC}&=\Delta_y+\mu_{x,A}-\mu_{x,C}
 -\mu_{y,A}-\mu_{y,C}+\mu_{xy,A}+\mu_{xy,C},\nonumber\\
 v^+_{AC}&=\Delta_y-\mu_{x,A}+\mu_{x,C}
 -\mu_{y,A}-\mu_{y,C}-\mu_{xy,A}-\mu_{xy,C}.
 \label{eq:verticalfaces}
\end{align}

Let \(t_x,t_y\) be block signs assigned by \eqref{eq:momentsign} on each
coarse row and column, and let \(s_x,s_y\) be the corresponding coarse face
signs. The face-current admissible set is
\begin{equation}
\begin{aligned}
 t_x(\mu_x\pm\mu_{xy})&\geq0,&s_xh^\pm&\geq0,\\
 t_y(\mu_y\pm\mu_{xy})&\geq0,&s_yv^\pm&\geq0.
\end{aligned}
 \label{eq:facepolytope}
\end{equation}
This admissible set is a finite intersection of affine half-spaces containing
the zero moment field. Admission first enforces the internal cone
\begin{equation}
 t_x\mu_x\geq|\mu_{xy}|,\qquad
 t_y\mu_y\geq|\mu_{xy}|.
 \label{eq:internalcone}
\end{equation}
For a cross face written \(g=\Delta+r_A+r_B\), only opposed proposal mass
consumes the witnessed capacity. Define
\begin{align}
 w_A&=\max(-sr_A,0),&w_B&=\max(-sr_B,0),\nonumber\\
 \beta&=\min\left(1,\frac{|\Delta|}{w_A+w_B}\right),
 \label{eq:facecapacity}
\end{align}
with \(\beta=1\) when the denominator is zero. Both incident block scales are
bounded by \(\beta\), and each block retains the minimum bound from all of its
faces. Since helpful components are not required for feasibility, a later
reduction on another face cannot violate an admitted inequality.
Consequently, \eqref{eq:internalcone}--\eqref{eq:facecapacity} yield a finite
local box inside \eqref{eq:facepolytope}.

Every pair of fine rows has a horizontal sign sequence interleaving
\(t_x,s_x,t_x,s_x,\ldots\); every pair of fine columns has the analogous
vertical sequence. The one-dimensional sign-count proof applies to each.
The two-dimensional restriction response is
\begin{equation}
 H(\omega_x,\omega_y)=e^{-\mathrm i(\omega_x+\omega_y)/2}
 \cos(\omega_x/2)\cos(\omega_y/2),
 \label{eq:nyquist2d}
\end{equation}
which annihilates both axial Nyquist lines and their checkerboard
intersection. Exact reconstruction is then obtained by restoring the three
residual moment fields in \eqref{eq:blockmoments}.

\subsection{Multilevel state}

Repeated analysis produces the state
\begin{equation}
 \mathcal L^{(J)}x=
 \left(z^{(J)},r^{(J)},r^{(J-1)},\ldots,r^{(1)}\right).
 \label{eq:multilevel}
\end{equation}
In one dimension, \(r^{(j)}\) is a centered-moment field; in two dimensions,
it contains the horizontal, vertical, and mixed residual moments. Analysis
continues only while every resulting coarse axis has at least five cells, and
synthesis restores the stored levels in reverse order. Induction on
\eqref{eq:fullcycle} then gives
\begin{equation}
 (\mathcal L^{(J)})^{-1}\mathcal L^{(J)}=I.
 \label{eq:multilevelinverse}
\end{equation}
For a one-dimensional fine line of length \(2N\), the \(N\) parent averages
and \(N\) residual moments are necessary and sufficient coordinates. For a
two-dimensional \(2N\times2M\) raster, the \(NM\) means and three \(NM\)
residual moment fields likewise account for every scalar degree of freedom.
